 \section{Introdução}

A deterioração clínica de pacientes hospitalizados representa um importante desafio para a garantia da saúde, especialmente em enfermarias, onde o monitoramento é realizado de forma intermitente \cite{vincent2018deterioracao}. Esse quadro se caracteriza pela piora progressiva de parâmetros fisiológicos, podendo evoluir para desfechos graves, como uma parada cardiorrespiratória, uma transferência não planejada para a Unidade de Terapia Intensiva (UTI) e até um óbito hospitalar \cite{gerry2020ews}. Portanto, é de suma importância propor métodos eficazes e de baixo custo para uma detecção ágil dessas descompensações, contribuindo assim para a redução de seus impactos.

Nesse cenário, destacam-se os Sistemas de Alerta Precoce (EWS, do inglês \textit{Early Warning Scores}), que consistem em ferramentas padronizadas de avaliação clínica \cite{ALAM2014587}. Esses sistemas têm como objetivo auxiliar na identificação antecipada da deterioração por meio da análise sistematizada de parâmetros fisiológicos rotineiramente coletados durante os cuidados ao paciente. A partir da atribuição de pontuações a variáveis como frequência respiratória, pressão arterial, frequência cardíaca, temperatura e saturação de oxigênio, obtém-se um escore capaz de indicar o grau de risco do paciente. Entre as implementações mais consagradas dessa metodologia está a Pontuação Nacional de Alerta Precoce (NEWS, do inglês \textit{National Early Warning Score}) e sua versão atualizada, o NEWS2 \cite{gerry2020ews}. Estudos apontam que o uso de ferramentas baseadas em EWS apresenta bom desempenho na identificação de eventos críticos iminentes em janelas de curto prazo, como nas primeiras 48 horas \cite{smith2014ews}.

No entanto, segundo \citeonline{scott2026news}, apesar de sua eficácia comprovada, a aplicação tradicional desses escores enfrenta dificuldades operacionais críticas. Ela depende de cálculos realizados manualmente à beira do leito pela equipe de enfermagem, utilizando-se de dados obtidos de forma pontual e espaçada. Esse modelo de aferição inconstante cria extensas janelas temporais em que o paciente permanece sem qualquer monitoramento, o que atrasa o reconhecimento do sintoma fisiológico e representa um obstáculo severo em situações de emergência.

Para superar essas barreiras físicas e operacionais, o desenvolvimento de sistemas embarcados apresenta-se como uma solução robusta, uma vez que a personalização do sistema embarcado permite adequá-lo de forma mais eficiente às necessidades apresentadas \cite{embedesSystem}. Dessa maneira, o equipamento pode coletar os sinais vitais continuamente e processá-los localmente, sem a necessidade de enviar dados brutos de telemetria de forma ininterrupta, o que previne problemas associados à infraestrutura operacional da rede hospitalar \cite{rancea2024edge}.

% Para superar essas barreiras físicas e operacionais, o desenvolvimento de sistemas embarcados se apresenta como uma solução robusta \textcolor{red}{Referência}. \textcolor{blue}{A construção de um equipamento dedicado, operando sob um sistema operacional Linux embarcado e customizado através de ferramentas de compilação como o \textit{Buildroot}, permite extrair o máximo de performance de um hardware e alcançar o baixo custo \cite{pera2022buildroot}. Essa abordagem garante a estabilidade, a segurança e o determinismo necessários para aplicações médicas \cite{l4b2021linux}, viabilizando o processamento de dados na borda (\textit{Edge Computing}).} Dessa forma, o equipamento pode coletar os sinais vitais continuamente e processá-los localmente, sem a necessidade de enviar dados brutos de telemetria de forma ininterrupta, o que previne a sobrecarga da infraestrutura de rede hospitalar \cite{rancea2024edge}.

Integrado a esse ecossistema de hardware restrito, o uso de algoritmos de Aprendizagem de Máquina (ML, do inglês \textit{Machine Learning}) tem se mostrado cada vez mais relevante na estratificação de risco clínico. Sendo um subcampo da Inteligência Artificial, essa tecnologia permite aos sistemas inferirem resultados complexos a partir de conjuntos de dados \cite{ibmMl}. Os modelos de ML podem ser divididos em supervisionados e não supervisionados. Em modelos supervisionados, foco deste trabalho, o sistema usa conjuntos de dados rotulados para entender as relações entre entradas e dados de saída, para poder trabalhar na previsão de saídas de novos conjuntos de dados \cite{ibmSupervised}.

Diante desse contexto, o presente trabalho tem como objetivo o desenvolvimento de um sistema embarcado para a automatização e predição contínua do protocolo EWS. A proposta se fundamenta na necessidade de conferir maior eficiência, padronização e escalabilidade ao processo, reduzindo a dependência de aferições intermitentes e mitigando falhas humanas. Assim, pretende-se oferecer uma solução de monitoramento constante e inteligente que apoie o reconhecimento precoce do agravamento clínico, favorecendo a tomada de decisão em tempo oportuno e elevando a segurança assistencial nas enfermarias.

A seguir, o artigo está organizado da seguinte forma: na Seção \ref{fundamentao}, são apresentados os fundamentos teóricos que embasam o trabalho proposto, com destaque em abordagens computacionais no contexto de Embarcados e Aprendizagem de Máquina. A Seção \ref{Correlatos} analisa os trabalhos correlatos, enquanto a Seção \ref{Metodologia} descreve a metodologia adotada. Na Seção \ref{Resultados}, discutem-se os experimentos realizados e os resultados obtidos.